%% ma: L12-B08-0182
%% lop: 12
%% chuong: 2
%% bai: 8
%% dang: 12.8.4
%% loai: TN4
%% muc: VD
%% dap_an: "A"
%% nguon: "CĐ2.6.pdf giáo viên gửi 27/10/2026 (M07, Lê Minh Khôi – Nguyễn Anh Khôi, Chương II Chủ đề 6 Ứng dụng vào các bài toán thực tế), Câu 5"
%% kien_thuc: "Bài 8 (tọa độ vectơ, độ dài), chuyển động thẳng đều (vectơ vận tốc)"
%% trang_thai: chinh-thuc
\begin{ex}
Trong không gian $Oxyz$ cho trước (đơn vị dài trên các trục là ki-lô-mét), trong khoảng thời gian $10$ phút, ra đa phát hiện hai máy bay chiến đấu di chuyển với vận tốc và hướng không đổi: máy bay (I) bay từ điểm $A(100;50;10)$ đến điểm $B(240;100;12)$; máy bay (II) bay từ điểm $C(80;40;6)$ đến điểm $D(220;100;9)$. Nếu hai máy bay giữ nguyên vận tốc và hướng bay thì khoảng cách giữa hai máy bay sau $15$ phút tiếp theo gần nhất với giá trị nào sau đây?
\choice{\True $25{,}04$ km}{$30{,}12$ km}{$20{,}50$ km}{$28{,}35$ km}
\loigiai{Sau $15$ phút: máy bay (I) ở $B+1{,}5\overrightarrow{AB}=(450;175;15)$, máy bay (II) ở $D+1{,}5\overrightarrow{CD}=(430;190;13{,}5)$.

Khoảng cách $\sqrt{20^2+15^2+1{,}5^2}=\sqrt{627{,}25}\approx25{,}04$ km.}
\end{ex}
